\chapter{Charged Particles, Gauge Invariance, and Poisson Brackets}

These companion notes follow Leonard Susskind's lecture closely, with the included board frames preserved as documentary evidence from the lecture materials curated by LazyingArt LLC and Video2Book. The lecture is not trying to teach electromagnetism in general. It returns to one example, a charged particle in prescribed electromagnetic fields, because this example exposes something peculiar about the Hamiltonian language: canonical momentum becomes surprising, gauge redundancy becomes unavoidable, and when an electric field is added the particle responds in the transverse way we call the Hall effect. The chapter should therefore read as a guided unfolding of that example, not as an anonymous summary of standard facts.

\section{Why Revisit the Charged Particle?}

We begin exactly where the lecture begins: the electromagnetic field itself is not being treated as dynamical. The electric and magnetic fields are simply given to us. For simplicity they are taken to be independent of time, though they may vary from point to point in space. Against that fixed background we watch a single charged particle move.

Why return to a problem we have already seen? Because, as Susskind says, the example is unintuitive. In a magnetic field the canonical formulation is not what one would have guessed in advance, the notion of momentum shifts under our feet, and once an electric field is added there is a surprising sideways response. So the lecture deliberately backs up and lists the hierarchy of formulations we have accumulated:
\begin{enumerate}
\item equations of motion,
\item the principle of least action,
\item the Lagrange equations,
\item the Hamiltonian formulation.
\end{enumerate}

That list is not decorative. The lecture uses it to rank what each form reveals. Equations of motion are the lowest common denominator: every system has them, but by themselves they do not display the deeper special structure of the laws. The principle of least action underlies everything, but in that raw form it is not yet a practical tool. The Lagrangian form is extremely powerful and general. The Hamiltonian form is in some ways the least intuitive, but it is the one that shows the flow in phase space and points directly toward quantum mechanics. The charged particle is the place where this whole ladder becomes concrete.

\section{From \(F=ma\) to the Planar Lorentz Force}

The lecture now narrows to a particle moving in two dimensions, simply because that is easy to visualize on the board. The magnetic field points perpendicular to the plane, into the board, and the electric field, when we restore it later, will lie in the plane. The nonrelativistic equation of motion is the Lorentz-force law,
\begin{equation}
m\ddot{\mathbf x}=q\mathbf E+q\,\dot{\mathbf x}\times\mathbf B.
\end{equation}
The lecture suppresses the factor of \(c\) that appears in some unit conventions; the structure of the equation is what matters here.

In the planar setup, \(\dot{\mathbf x}\) lies in the plane and \(\mathbf B\) points out of it, so \(\dot{\mathbf x}\times\mathbf B\) again lies in the plane. Thus the full motion closes within the plane. This is why the two-dimensional example is not a loss of substance but a simplification of geometry.

The first physical fact extracted from the equation is immediate and important. The magnetic force depends on the velocity, but it is perpendicular to the velocity. Therefore
\begin{equation}
\mathbf v\cdot(\mathbf v\times\mathbf B)=0.
\end{equation}
This is the concise form of the lecture's slogan: magnetic fields do no work. They deflect the trajectory, but they do not change the speed and therefore do not change the kinetic energy.

The electric force is of a more ordinary sort. In the static situation under discussion it is conservative, and so it belongs to the familiar language of potential energy:
\begin{equation}
\mathbf F=-\nabla U.
\end{equation}

\begin{figure}[t]
\centering
\includegraphics[width=0.56\textwidth]{lecture_08_figure_03.png}
\caption{Conservative force written as a gradient.}
\end{figure}

This is a small but important pivot in the lecture. The electric field can be absorbed into the ordinary machinery of conservative mechanics. The magnetic field cannot. That is why the lecture now interrupts the forward motion and introduces a new object.

\section{Potentials and the Action}

Before we can move from the force law to least action, we need one more idea: the vector potential. The magnetic field is divergence-free, and the lecture pauses to recall what that means physically. At least in the physics ordinarily done in this regime, there are no magnetic charges serving as sources of \(\mathbf B\). A divergence-free field can be expressed as the curl of something else, and so we write
\begin{equation}
\mathbf B=\nabla\times\mathbf A.
\end{equation}
But not uniquely. The vector potential is ambiguous from the start, and the lecture flags that ambiguity before it is used.

The electric field, in contrast, is introduced in the more familiar conservative form. Since the force on a charge is \(q\mathbf E\), one writes
\begin{equation}
\mathbf E=-\nabla V,
\qquad
U=qV.
\end{equation}
The scalar \(V\) is the electric potential, an energy per unit charge.

\begin{figure}[t]
\centering
\includegraphics[width=0.58\textwidth]{lecture_08_figure_04.png}
\caption{Vector potential, electric potential, and \(U=qV\).}
\end{figure}

The board layout here is worth preserving because it matches the lecture's logic: first \(\mathbf B\) from \(\mathbf A\), then \(\mathbf E\) from \(V\), then \(U=qV\). Only after that chain is in place does the action appear.

For motion in the \(x\)-\(y\) plane, the action is written as
\begin{equation}
S=\int \frac12 m(\dot x^2+\dot y^2)\,dt+q\int A_i\,dx_i-q\int V\,dt.
\end{equation}
The first term is the usual kinetic term. The last term is the scalar-potential contribution. The middle term is the genuinely new ingredient: the magnetic field enters through a line integral of the vector potential along the trajectory. The board shot is slightly abbreviated in its last term, but once the lecture has established \(U=qV\), the clean formula must carry the factor of \(q\) in front of \(V\).

\begin{figure}[t]
\centering
\includegraphics[width=0.78\textwidth]{lecture_08_figure_05.png}
\caption{Electromagnetic action arranged as a Lagrangian integral.}
\end{figure}

At first this does not look like the sort of action we have been using, because \(A_i\,dx_i\) is written in terms of infinitesimal displacements rather than velocities. But the lecture immediately massages it into ordinary Lagrangian form by dividing and multiplying by \(dt\):
\begin{equation}
A_i\,dx_i=A_i\,\dot x_i\,dt.
\end{equation}
Hence
\begin{equation}
S=\int\left[\frac12 m(\dot x^2+\dot y^2)+qA_i\dot x_i-qV\right]dt,
\end{equation}
so the Lagrangian is
\begin{equation}
L=\frac12 m(\dot x^2+\dot y^2)+qA_i\dot x_i-qV
=\frac12 m(\dot x^2+\dot y^2)+q\,\mathbf A\cdot\mathbf v-qV.
\end{equation}

The lecture lingers for a moment on the form of this expression. The spatial terms \(A_i\,dx_i\) and the temporal term \(V\,dt\) are beginning to arrange themselves as though they belonged to a larger space-time object. He does not pursue relativity here, but he does let us see the direction in which the formula is leaning.

\subsection{Question \& Answer}

Question. Why must the action use \(\mathbf A\) if the motion depends only on \(\mathbf B\)?

Answer. This is the lecture's first real conceptual tension, and Susskind deliberately leaves it visible. The equations of motion depend on the magnetic field, not on the vector potential. One can change the vector potential without changing \(\mathbf B\), and the particle's motion does not notice. Nevertheless, the only known action that reproduces the Lorentz force requires \(\mathbf A\). We are forced to write down more than the observable motion seems to need.

The resolution is gauge invariance. We may transform
\begin{equation}
\mathbf A\to\mathbf A+\nabla\lambda,
\end{equation}
and because
\begin{equation}
\nabla\times(\nabla\lambda)=0,
\end{equation}
the magnetic field is unchanged. In the static setting of the lecture, this changes the Lagrangian by
\begin{equation}
q\,\nabla\lambda\cdot\mathbf v
=
q\,\frac{d\lambda}{dt},
\end{equation}
a total derivative along the trajectory. So the Lagrangian changes, but the Euler-Lagrange equations do not. The extra description is redundant, not additional physics. That redundancy is precisely what the lecture means by gauge invariance.

\section{Deriving the Lorentz Force and Exposing Canonical Momentum}

The lecture now does something important stylistically as well as mathematically: it refuses to leave the formalism at the level of assertion. Instead it derives the Lorentz force back from the Lagrangian, and it does so the old-fashioned way, by taking a single component and working through it carefully. Let
\begin{equation}
L=\frac12 m(\dot x^2+\dot y^2)+q(A_x\dot x+A_y\dot y)-qV(x,y).
\end{equation}
We do only the \(x\)-component, just as in the lecture.

First compute the canonical momentum conjugate to \(x\):
\begin{equation}
p_x=\frac{\partial L}{\partial \dot x}=m\dot x+qA_x.
\end{equation}
This is already the surprise the lecture wanted us to notice. Canonical momentum is not just \(m\dot x\). There is an extra term, and that extra term depends on the vector potential.

\begin{figure}[t]
\centering
\includegraphics[width=0.78\textwidth]{lecture_08_figure_06.png}
\caption{Pointing to the \(\dot x\)-dependent term in the Lagrangian.}
\end{figure}

Now apply the Euler-Lagrange equation,
\begin{equation}
\frac{d}{dt}\frac{\partial L}{\partial \dot x}=\frac{\partial L}{\partial x}.
\end{equation}
The left-hand side is
\begin{equation}
\frac{d}{dt}\frac{\partial L}{\partial \dot x}
=
m\ddot x+q\frac{dA_x}{dt}.
\end{equation}
Here the lecture makes a useful clarification. The field is assumed not to depend explicitly on time, but that does not mean \(A_x\) is constant along the particle's trajectory. As the particle moves through space, the value of the field at the particle changes. Therefore
\begin{equation}
\frac{dA_x}{dt}
=
\frac{\partial A_x}{\partial x}\dot x
+
\frac{\partial A_x}{\partial y}\dot y.
\end{equation}

On the right-hand side,
\begin{equation}
\frac{\partial L}{\partial x}
=
q\frac{\partial A_x}{\partial x}\dot x
+
q\frac{\partial A_y}{\partial x}\dot y
-q\frac{\partial V}{\partial x}.
\end{equation}
Putting the two sides together,
\begin{align}
m\ddot x
+q\frac{\partial A_x}{\partial x}\dot x
+q\frac{\partial A_x}{\partial y}\dot y
&=
q\frac{\partial A_x}{\partial x}\dot x
+
q\frac{\partial A_y}{\partial x}\dot y
-q\frac{\partial V}{\partial x}.
\end{align}
The \(\partial_xA_x\,\dot x\) terms cancel, leaving
\begin{align}
m\ddot x
&=
q\left(\frac{\partial A_y}{\partial x}-\frac{\partial A_x}{\partial y}\right)\dot y
-q\frac{\partial V}{\partial x} \\
&= qB_z\,\dot y+qE_x,
\end{align}
where
\begin{equation}
B_z=\frac{\partial A_y}{\partial x}-\frac{\partial A_x}{\partial y},
\qquad
E_x=-\frac{\partial V}{\partial x}.
\end{equation}
This is exactly the \(x\)-component of the Lorentz force. The calculation earns its keep in two ways at once: it verifies the Lagrangian, and it shows that the vector potential reenters the equations only through the magnetic field itself.

The lecture then turns back to the momentum formula and lets the oddity stand. If
\begin{equation}
p_x=m\dot x+qA_x,
\end{equation}
then \(p_x\) changes under a gauge transformation. That means canonical momentum is not gauge invariant. Susskind's attitude here is direct: if the mathematics tells us that the canonical momentum in a magnetic field is gauge-dependent, then we follow the mathematics. The real question is what physical role this quantity is actually playing.

\section{Uniform Magnetic Field and the Meaning of the Conserved Momentum}

To answer that question, the lecture turns from derivation to exploitation. We now take a uniform magnetic field \(B_z=B\), constant in space. Several gauges are written on the board. Their signs wobble a little in the transcript, but the later conservation-law argument uses the pair
\begin{equation}
A_x=By,\qquad A_y=0,
\end{equation}
and
\begin{equation}
A_x=0,\qquad A_y=-Bx.
\end{equation}
These are two different gauges for the same field.

Take the first gauge and, for the moment, set the electric field to zero. Then
\begin{equation}
L=\frac12 m(\dot x^2+\dot y^2)+qBy\,\dot x.
\end{equation}
The canonical momenta are
\begin{equation}
p_x=m\dot x+qBy,
\qquad
p_y=m\dot y.
\end{equation}
Now the lecture asks a symmetry question rather than solving an equation. The Lagrangian depends on \(y\), but not on \(x\). So a translation \(x\to x+\epsilon\) leaves \(L\) unchanged, while a translation \(y\to y+\epsilon\) does not. Therefore \(p_x\) is conserved, while \(p_y\) is not. In the especially simple case \(p_x=0\), we learn at once that
\begin{equation}
\dot x=-\frac{qB}{m}\,y.
\end{equation}

Now change gauge. In the second gauge,
\begin{equation}
L=\frac12 m(\dot x^2+\dot y^2)-qBx\,\dot y,
\end{equation}
and so
\begin{equation}
p_y=m\dot y-qBx.
\end{equation}
This time the Lagrangian depends on \(x\), but not on \(y\). So now \(p_y\) is conserved, and in the simple case \(p_y=0\),
\begin{equation}
\dot y=\frac{qB}{m}\,x.
\end{equation}

This is the lecture's peculiar but perfectly sound logic: one gauge gives us the \(x\)-piece of the story, another gives us the \(y\)-piece, and the physical motion is entitled to both because the motion itself does not depend on the gauge.

\subsection{Question \& Answer}

Question. What is the conserved momentum actually measuring in a uniform magnetic field?

Answer. The first step is to combine the two gauge-invariant velocity relations:
\begin{equation}
\dot x=-\frac{qB}{m}\,y,
\qquad
\dot y=\frac{qB}{m}\,x.
\end{equation}
These are precisely the relations one expects for uniform circular motion. To make that explicit, take
\begin{equation}
x=r\cos(\omega t),
\qquad
y=r\sin(\omega t).
\end{equation}
Then
\begin{equation}
\dot x=-r\omega\sin(\omega t)=-\omega y,
\qquad
\dot y=r\omega\cos(\omega t)=\omega x.
\end{equation}
Comparison gives the cyclotron frequency
\begin{equation}
\omega=-\frac{qB}{m},
\end{equation}
up to the sign convention fixed by the orientation of the magnetic field and the sign of the charge. As the lecture notes, experiments of this kind determine the ratio \(q/m\), not \(q\) and \(m\) separately.

At this point the lecture asks us to push one step further. A uniform magnetic field does not distinguish the origin, so why should the circle be centered there? Let us shift the orbit:
\begin{equation}
x=x_0+r\cos(\omega t),
\qquad
y=y_0+r\sin(\omega t).
\end{equation}
The velocities are unchanged by the shift. Now substitute into the conserved quantities. Using \(\omega=-qB/m\), the time-dependent terms cancel, and one finds
\begin{equation}
p_x=qB\,y_0,
\qquad
p_y=-qB\,x_0.
\end{equation}
This is the real meaning of the conserved canonical momenta. They are not encoding translational motion in the naive sense. They are encoding the fixed guiding-center coordinates of the cyclotron orbit. Momentum conservation has been reorganized into a statement that the center of the circle does not move.

This is exactly the sort of thing the lecture wanted the example to expose. The canonical formalism is strange here, but it is not wrong. On the contrary, once we accept the gauge redundancy, the conservation laws extract the motion in an almost embarrassingly efficient way.

\section{Electric Field, Drift, and the Hall Effect}

Only after the pure magnetic problem has been fully digested does the lecture put the electric field back in. Take the electric field to point in the \(+x\)-direction. To be consistent with
\begin{equation}
\mathbf E=-\nabla V,
\end{equation}
we choose
\begin{equation}
V(x)=-Ex,
\end{equation}
so that
\begin{equation}
\mathbf E=E\,\hat{\mathbf x}.
\end{equation}
For the same uniform magnetic field as before, the equations become
\begin{align}
m\ddot x+qB\,\dot y &= qE, \\
m\ddot y-qB\,\dot x &= 0.
\end{align}
The first equation records that the \(x\)-translation symmetry has been broken; \(p_x\) is no longer conserved. The second remains unchanged because no force has been added in the \(y\)-direction.

But the lecture does not solve the full system. It asks for a more special question: is there a solution with no acceleration at all?

\subsection{Question \& Answer}

Question. How can an electric force in the \(x\)-direction produce motion in the \(y\)-direction?

Answer. Assume
\begin{equation}
\ddot x=\ddot y=0.
\end{equation}
Then the second equation gives
\begin{equation}
\dot x=0,
\end{equation}
and the first reduces to
\begin{equation}
qB\,\dot y=qE.
\end{equation}
Thus
\begin{equation}
\dot y=\frac{E}{B},
\end{equation}
for the orientation convention used in the lecture. Reversing the orientation of \(\mathbf B\) reverses the drift, but the scale is always set by the ratio \(E/B\).

This is the payoff that was promised at the beginning. We applied a force along \(x\), but the distinguished no-acceleration response is motion along \(y\). That is the Hall effect in this stripped-down mechanical setting. The lecture sharpens the surprise further by observing that without the electric field the only no-acceleration solution is the trivial one, the zero-radius circle in which the particle simply sits still. Once the electric field is turned on, the particle is pushed into steady transverse motion.

There is another point the lecture emphasizes. The drift speed does not depend on the magnitude of the charge; the \(q\)'s cancel. So if a whole distribution of charge sits in the plane under a uniform magnetic field, turning on an electric field along \(x\) drives the whole distribution sideways along \(y\), with speed determined by \(E/B\). Susskind remarks that the mathematics is close to gyroscopic precession: an applied influence in one direction produces motion in the perpendicular direction.

\section{Poisson Brackets as Another Form of Mechanics}

The lecture could have continued with the charged particle, perhaps by constructing the Hamiltonian explicitly. Instead Susskind leaves that as a problem and deliberately changes register. We have, as he says, been through electric particles, magnetic particles, and gauge invariance. Now he wants to come to another form of mechanics, very closely related to Hamilton's form but more abstract: the form in terms of Poisson brackets.

He is explicit about the motivation. This is not being introduced because it is the most intuitive classical language. It is being introduced because it is central to quantum mechanics, and because at some point one begins to suspect that classical mechanics itself is already a disguised statement about the motion of quantum wave packets.

Start with Hamilton's equations,
\begin{equation}
\dot q_i=\frac{\partial H}{\partial p_i},
\qquad
\dot p_i=-\frac{\partial H}{\partial q_i}.
\end{equation}
Let \(A(q_i,p_i)\) be any function on phase space, with no explicit time dependence. Along a trajectory,
\begin{align}
\dot A
&=
\sum_i\frac{\partial A}{\partial q_i}\dot q_i
+\sum_i\frac{\partial A}{\partial p_i}\dot p_i \\
&=
\sum_i
\left(
\frac{\partial A}{\partial q_i}\frac{\partial H}{\partial p_i}
-\frac{\partial A}{\partial p_i}\frac{\partial H}{\partial q_i}
\right).
\end{align}
This combination is elevated to a definition. For any two phase-space functions \(A\) and \(B\),
\begin{equation}
\{A,B\}
=
\sum_i
\left(
\frac{\partial A}{\partial q_i}\frac{\partial B}{\partial p_i}
-
\frac{\partial A}{\partial p_i}\frac{\partial B}{\partial q_i}
\right).
\end{equation}
Then time evolution becomes
\begin{equation}
\dot A=\{A,H\}.
\end{equation}

That single formula already contains Hamilton's equations. For example, in one degree of freedom,
\begin{equation}
\dot Q=\{Q,H\}=\frac{\partial H}{\partial P},
\qquad
\dot P=\{P,H\}=-\frac{\partial H}{\partial Q}.
\end{equation}
So the bracket does not replace Hamiltonian mechanics; it compresses it into an algebraic rule.

The lecture then builds the algebra from a few basic identities. First,
\begin{equation}
\{q_i,p_j\}=\delta_{ij},
\qquad
\{q_i,q_j\}=0,
\qquad
\{p_i,p_j\}=0.
\end{equation}
The bracket is antisymmetric,
\begin{equation}
\{A,B\}=-\{B,A\},
\end{equation}
so in particular
\begin{equation}
\{p_i,q_j\}=-\delta_{ij}.
\end{equation}

For a single canonical pair \((Q,P)\), the bracket acts like differentiation:
\begin{equation}
\{P,F\}=-\frac{\partial F}{\partial Q},
\qquad
\{Q,F\}=\frac{\partial F}{\partial P}.
\end{equation}
That is, bracketing with \(P\) differentiates with respect to \(Q\), while bracketing with \(Q\) differentiates with respect to \(P\).

One also needs linearity,
\begin{equation}
\{\alpha A,B\}=\alpha\{A,B\},
\qquad
\{A+C,B\}=\{A,B\}+\{C,B\},
\end{equation}
and the product rule,
\begin{equation}
\{AB,C\}=A\{B,C\}+B\{A,C\}.
\end{equation}
With the fundamental brackets, antisymmetry, linearity, the product rule, and the rule \(\dot A=\{A,H\}\), the lecture says one has, in effect, a compact algebraic formulation of classical mechanics itself.

The charged particle has now done all the work it needed to do. It revealed where the Hamiltonian language becomes unintuitive, and then the lecture pivoted to a formulation in which the underlying structure is stripped down to brackets and rules. The point of that pivot is not local elegance. It is preparation for the next encounter with quantum mechanics.

\section{Summary}

The lecture began by returning to a charged particle in prescribed electromagnetic fields because the example is richer than it first appears. Starting from the Lorentz-force law, we saw that the magnetic field bends motion without doing work. Passing to least action forced us to introduce the vector potential, and that in turn forced us to confront gauge redundancy: the action needs \(\mathbf A\), even though the observable motion depends only on \(\mathbf B\). Once the Lagrangian is written down, the canonical momentum appears in a gauge-dependent form. In a uniform magnetic field, the conserved canonical quantities do not describe naive translational momentum but the fixed center of the cyclotron orbit. Restoring the electric field then produces the Hall drift, a transverse response with speed set by \(E/B\). Only after that physical story is complete does the lecture step back and reformulate Hamiltonian mechanics in terms of Poisson brackets, making the algebraic structure that leads toward quantum mechanics unmistakable.